% TOP_PROBLEM: {"id":30007724,"problem_number":"LOCAL-30007724","title":"Insurance choice frictions and selection","category_id":21,"category":{"id":21,"name":"economics","display_name":"Economics","description":"Open economic theory statements, published research agendas, and proposed scoped specifications; question provenance and historical priority are qualified per record.","slug":"economics","order_index":21,"created_at":"2026-10-08T00:00:00Z"},"status":"research_question","external_url":null,"legacy_ids":[],"published":false,"statement":"Measure whether an insurance decision aid improves welfare after risk sorting and insurers’ subsequent premium responses, rather than judge it by switching rates alone.","economics":{"original_id":213,"field":"Health economics","kind":"design","formalization_basis":"adapted_research_question","source_status":"research_agenda","priority_status":"earliest_found","priority_note":"Earliest explicit statement located in this audit; no exhaustive priority search or first-ever claim. The mathematical specification below is adapted, not quoted from the source.","earliest_verified_reference":{"authors":["Martin Gaynor","Kate Ho","Robert J. Town"],"title":"The Industrial Organization of Health Care Markets","year":2014,"url":"https://kateho.scholar.princeton.edu/sites/g/files/toruqf4136/files/kateho/files/healthcare_io_3_24_14.pdf","doi":"10.3386/w19800","locator":"§7 Models of Health Plan Choice, p. 40; §8, p. 43; §9 Future Directions, pp. 47–48","evidence_summary":"The authors explicitly discuss unanswered questions about consumer inertia, adverse selection, provider incentives, and consumer search frictions."},"current_reference":{"authors":["Martin Gaynor","Kate Ho","Robert J. Town"],"title":"The Industrial Organization of Health Care Markets","year":2014,"url":"https://kateho.scholar.princeton.edu/sites/g/files/toruqf4136/files/kateho/files/healthcare_io_3_24_14.pdf","doi":"10.3386/w19800","locator":"§7 Models of Health Plan Choice, p. 40; §8, p. 43; §9 Future Directions, pp. 47–48","evidence_summary":"The authors explicitly discuss unanswered questions about consumer inertia, adverse selection, provider incentives, and consumer search frictions."},"formalization_note":"The source explicitly identifies the underlying gap or agenda. Population, horizon, interventions, estimands, and auxiliary assumptions are proposed operational choices, not a canonical source theorem. Partial later answers are compatible with this scoped question; the citation alone does not prove absence of every subsequent solution."},"statusQualification":"Published question or research agenda verified at the cited locator. The mathematical specification is adapted for this catalogue and includes additional assumptions; source publication does not establish first-ever priority or certify this exact model remains unsolved. The source explicitly identifies the underlying gap or agenda. Population, horizon, interventions, estimands, and auxiliary assumptions are proposed operational choices, not a canonical source theorem. Partial later answers are compatible with this scoped question; the citation alone does not prove absence of every subsequent solution.","statusReviewedAt":"2026-10-08"}
% FIRST_POSTED: 2026-10-08
% ENTRY_KIND: statement_only
% !TeX program = lualatex
\documentclass[11pt]{article}
\usepackage{fontspec}
\usepackage[a4paper,margin=25mm]{geometry}
\usepackage{amsmath,amssymb,mathtools}
\usepackage{xurl}
\usepackage[hidelinks]{hyperref}
\newcommand{\sref}[2]{\hyperlink{source:#1:#2}{[S#2]}}
\setlength{\emergencystretch}{3em}
\begin{document}
% problemId:30007724
\section*{Insurance choice frictions and selection}\label{30007724}
\subsection{Definitions and mathematical statement}
Consider adults purchasing individual health insurance in one United States state marketplace over three enrollment years. Let \(z\in\{0,1\}\) denote randomized access to a specified decision aid that explains premiums, provider networks, and expected out-of-pocket expenditure. Insurers choose subsequent premiums in response to enrollment, under unchanged risk-adjustment rules. Define household certainty-equivalent utility \(U_i(z)\), measured in annual dollars after medical expenditure risk and valued provider access, and insurer profit \(\Pi_j(z)\). The equilibrium welfare contrast is
\[\Delta W=E\!\left[\sum_i U_i(1)+\sum_j\Pi_j(1)-\sum_iU_i(0)-\sum_j\Pi_j(0)\right]-K,\]
where \(K\) is the aid’s real implementation cost and transfers are not counted twice. Determine how reducing choice frictions changes welfare once risk sorting and premiums adjust, rather than evaluate plan-switching rates alone. Specify a model connecting measured preferences and risk to insurer pricing, and identify decision-aid effects through randomized market-level saturation or justified equilibrium restrictions. Distinguish observed inertia from a preference for continuity or switching costs. The older health-market review explicitly links unanswered inertia questions to adverse selection; this three-year marketplace experiment is an adapted research design.

\par\textbf{Formulation and scope.} The source explicitly identifies the underlying gap or agenda. Population, horizon, interventions, estimands, and auxiliary assumptions are proposed operational choices, not a canonical source theorem. Partial later answers are compatible with this scoped question; the citation alone does not prove absence of every subsequent solution.

\par\textbf{Open-question provenance.} An explicit question or research agenda was verified at the source locator. This does not by itself certify that every later version remains unresolved \sref{30007724}{1}.

\par\textbf{Historical priority.} Earliest explicit statement located in this audit; no exhaustive priority search or first-ever claim. The mathematical specification below is adapted, not quoted from the source.

\subsection{Short English statement}
Measure whether an insurance decision aid improves welfare after risk sorting and insurers’ subsequent premium responses, rather than judge it by switching rates alone.

\subsection{Sources}
\begin{itemize}\raggedright
\item[\textnormal{[S1]}]\hypertarget{source:30007724:1}{} Martin Gaynor, Kate Ho, Robert J. Town. 2014. The Industrial Organization of Health Care Markets. §7 Models of Health Plan Choice, p. 40; §8, p. 43; §9 Future Directions, pp. 47–48.
\url{https://kateho.scholar.princeton.edu/sites/g/files/toruqf4136/files/kateho/files/healthcare_io_3_24_14.pdf}.
DOI: 10.3386/w19800.
{\small\textit{Earliest verified question/agenda reference; historical priority qualified below. The authors explicitly discuss unanswered questions about consumer inertia, adverse selection, provider incentives, and consumer search frictions.}}
\end{itemize}
\end{document}
